\chapter{Скрученный спектр полного тензорного вихря на стационарном фоне}
\label{ch:v6-bosonic-defect-full-tensor-stationary-twisted-spectrum-gate}

\section{Постановка}

Четырёхпрофильный фон $(K,a,b,q)$ устранил ненулевую первую вариацию поля
$Q$. Поэтому полный полярный оператор теперь можно вычислять в физически
допустимой стационарной точке. Как и в предыдущей главе, голономия $Z_3$
разбивает восемнадцать комплексных радиальных полей на три характера по шесть
полей в каждом.

В оператор включены все пять компонент $Q$, семь компонент $T$, радиальная и
угловая компоненты каждой из трёх направлений семейной связности, фоновая
калибровка, неабелев перекрёстный член кривизны и скрученные условия
регулярности. Методическая необходимость отделить переносные нули от
внутренних уровней соответствует стандартному спектральному анализу вихрей
\cite{stationary-spectrum-alonso,stationary-spectrum-forgacs}.

\section{Полное окно скрученных блоков}

На сетках $70$, $100$, $140$ узлов проверены три характера
$c=-1,0,1$ и метки
\begin{equation}
 -3\le j\le3.
\end{equation}
Два нетривиальных характера образуют сопряжённые пары. Максимальное различие
соответствующих собственных значений равно
\begin{equation}
 1.77\cdot10^{-11},
\end{equation}
а максимальный эрмитов остаток --- $2.21\cdot10^{-14}$.

Нижняя сопряжённая пара $(c,j)=(1,-1),(-1,1)$ принимает значения
\begin{equation}
 -0.01785481,\qquad -0.02094520,\qquad -0.02238156.
\end{equation}
Вторая пара $(1,-2),(-1,2)$ равна
\begin{equation}
 -0.00696658,qquad -0.00915328,qquad -0.01013267.
\end{equation}
Обе последовательности при сгущении уходят от нуля в отрицательную область.

\section{Целевое сгущение}

Критические блоки дополнительно вычислены на сетках $200$ и $280$ узлов.
Для первой пары получено
\begin{equation}
 -0.02313895,qquad -0.02349349,
\end{equation}
а для второй
\begin{equation}
 -0.01063964,qquad -0.01087481.
\end{equation}
Линейная экстраполяция по $1/(N-1)^2$ даёт пределы
\begin{equation}
 \lambda_{1,\infty}=-0.02386058,
 \qquad
 \lambda_{2,\infty}=-0.01112057.
\end{equation}
Сдвиги между сетками $200$ и $280$ равны соответственно
$3.55\cdot10^{-4}$ и $2.35\cdot10^{-4}$.

Отрицательный кандидат, таким образом, не исчез после стационаризации фона и
не стремится к нулю при сгущении.

\section{Почему сертификат ещё не закрыт}

Нейтральный блок $c=0$, $j=\pm1$ должен содержать точные переносные нули.
Численно его нижний уровень равен
\begin{equation}
 0.00996187, 0.00802057, 0.00715253, 0.00669883, 0.00648619
\end{equation}
на сетках $70$, $100$, $140$, $200$, $280$. Формальная экстраполяция даёт
$0.00626629$, а не нуль. Следовательно, текущая дискретизация всё ещё
содержит положительное решёточное поднятие геометрического переноса.

Модуль отрицательной первой пары примерно в $3.7$ раза больше этого
остаточного масштаба, поэтому она является сильным кандидатом на физическую
неустойчивость. Но один и тот же оператор обязан сначала воспроизвести
известный точный нуль. Без такой калибровки нельзя исключить недостающий
перекрёстный член или неверную реализацию условия в сердцевине.

\begin{nogocriterion}
Сходящаяся отрицательная последовательность не объявляется физической модой,
пока дискретизированный полный оператор не обращает ковариантную переносную
касательную в нуль с той же точностью. Знак отрицательного кандидата
фиксируется, но его физический статус остаётся открытым.
\end{nogocriterion}

\section{Следующий различающий тест}

Нулевая гипотеза следующего гейта: положительный переносный остаток вызван
единственным пропущенным или неверно нормированным членом дискретной второй
вариации; после его исправления отрицательная скрученная пара исчезнет.

Стоп-критерий: если оператор проходит точную ковариантную переносную
калибровку, а предел $\lambda_{1,\infty}<0$ сохраняется, прямая вихревая
нить является линейным седлом полного действия.

\begin{protocol}
\begin{enumerate}
 \item стационарный четырёхпрофильный фон: \textbf{использован};
 \item проверенное окно: \textbf{$c=-1,0,1$, $-3\le j\le3$};
 \item сопряжённость нетривиальных секторов: \textbf{прошла};
 \item нижний отрицательный кандидат: \textbf{$-0.02386058$};
 \item второй отрицательный кандидат: \textbf{$-0.01112057$};
 \item переносный остаток: \textbf{$+0.00626629$};
 \item физическая отрицательная мода: \textbf{ещё не сертифицирована};
 \item полная устойчивость нити: \textbf{не закрыта};
 \item следующий гейт: \textbf{точная переносная калибровка полного оператора}.
\end{enumerate}
\end{protocol}

Машинный сертификат сохранён в
\path{s2t_v6_bosonic_defect_full_tensor_stationary_twisted_spectrum_gate_results.json}.

\begin{thebibliography}{9}
\bibitem{stationary-spectrum-alonso}
J.~M.~Alonso-Izquierdo, W.~Garc\'ia Fuertes, J.~Mateos Guilarte,
\emph{Dissecting zero modes and bound states on BPS vortices in
Ginzburg--Landau superconductors}, arXiv:1602.09084.

\bibitem{stationary-spectrum-forgacs}
P.~Forg\'acs, S.~Reuillon, M.~S.~Volkov,
\emph{Stability analysis of twisted superconducting semilocal strings},
arXiv:0712.3589.
\end{thebibliography}